\chapter{Rationality conditions for structured argumentation}
\label{sec:p18-structured-rationality}

This chapter uses the corrected March 1, 2018 author version of Modgil
and Prakken's general account, together with its corrigendum
\cite{ModgilPrakken2013GeneralCorrected,ModgilPrakken2018Corrigendum}.
Its preference convention differs from the tutorial convention in
Chapter~\ref{sec:p18-structured-argumentation}. We define that difference
explicitly. We then establish a general consistency theorem with an
explicit local return condition. An exact countermodel shows why closure
under transposition and the finite-set reading of the displayed
reasonable-order conditions cannot replace that condition for the full
class of unrestricted trees. Quantification over indexed families with
repeated arguments is a distinct, stronger condition.

\section{Source conventions and their domains}

Use the language, two rule types, knowledge base, finite argument trees,
contrary function, attacks, and target-local defeats already defined.
For this source presentation, names may be supplied for every defeasible
rule; the results also allow the earlier partial naming convention.
Write \(T\) for structural attack and \(D\) for successful defeat.
For an argument \(A\), let \(\operatorname{Str}(A)\) be its set of strict
rules, and let \(\operatorname{Prem}_n(A)=\operatorname{Prem}(A)\cap K_n\).
Call an argument \emph{fallible} if it has an ordinary premise or a
defeasible rule. A strict firm argument is therefore precisely an
argument that is not fallible.

\begin{hbdefinition}{def:p18-rationality-support}{Two argument universes}
For \(S\subseteq L\), the notation \(S\vdash_s\phi\) means that a finite
strict tree derives \(\phi\) with premises in \(S\). In this consequence
notation, members of \(S\) may be used as premise leaves even when they
are not in the fixed knowledge base.
A set \(S\) is \emph{c-consistent} if its strict closure contains no
pair of contradictory formulas. One-way contraries do not by themselves
violate c-consistency.

An unrestricted structured framework uses all finite argument trees over
the fixed knowledge base. A c-restricted framework uses exactly the trees
whose premise sets are c-consistent. It imposes no minimal-support
condition. Subtrees of an admitted c-restricted tree are admitted, since
c-consistency is inherited by subsets.
\end{hbdefinition}

\begin{hbdefinition}{def:p18-rationality-closure}{Closure and background conditions}
Axiom consistency means that \(\operatorname{Cl}_s(K_n)\) is directly
consistent. Contrary well-formedness means that, whenever \(\alpha\) is
a one-way contrary of \(\beta\), the target \(\beta\) is neither an axiom
nor the consequent of a strict rule.

Closure under contraposition means that, for every \(S\subseteq L\),
\(s\in S\), and \(\phi\), if \(S\vdash_s\phi\), then
\[
 (S\setminus\{s\})\cup\{\phi^\ast\}\vdash_s s^\ast
\]
for every chosen contradictory \(\phi^\ast\) of \(\phi\), with a
contradictory \(s^\ast\) of \(s\) as consequent.
In a paired-literal language, closure under transposition is the
rule-by-rule condition of Definition~\ref{def:p18-aspic-consistency}.
When antecedents are ordered lists, include their permutations as well;
the finite examples below do so explicitly.

The c-classicality condition says that every inclusion-minimal
c-inconsistent set \(S\) satisfies
\(S\setminus\{\phi\}\vdash_s\phi^\ast\) for each \(\phi\in S\),
for some contradictory \(\phi^\ast\) of \(\phi\).
Only finite minimal sets are needed below: a c-inconsistency has two
finite derivations and therefore a finite witness.
\end{hbdefinition}

\begin{hbdefinition}{def:p18-rationality-reasonable}{The source's reasonable-order conditions}
The maximal fallible frontier \(M(A)\) consists of subtrees whose root
is an ordinary premise or a defeasible rule and which are not proper
subtrees of another such subtree of \(A\). Equivalently, traverse each
root-to-leaf path and stop at its first fallible node. If \(A\)'s root
is fallible, \(M(A)=\{A\}\).

The \emph{profile continuation} relation in the general account declares
\(B\) to be a continuation of a finite family \(\mathcal H\) when
\[
\begin{aligned}
 \operatorname{Prem}_p(B)&=\bigcup_{A\in\mathcal H}\operatorname{Prem}_p(A),
 &\operatorname{Def}(B)&=\bigcup_{A\in\mathcal H}\operatorname{Def}(A),\\
 \operatorname{Prem}_n(B)&\supseteq\bigcup_{A\in\mathcal H}\operatorname{Prem}_n(A),
 &\operatorname{Str}(B)&\supseteq\bigcup_{A\in\mathcal H}\operatorname{Str}(A).
\end{aligned}
\]
This is a condition on component sets. It does not say that the input
trees occur as actual subtrees of \(B\).

For an asymmetric strict preference \(\prec\), the source's
reasonable-order conditions are:

(i) every fallible argument is strictly below every strict firm argument;

(ii) a strict firm argument is never strictly below another argument;

(iii) if \(A'\) is a profile continuation of \(\{A\}\), then
\[
 A\not\prec B\Longrightarrow A'\not\prec B,\qquad
 B\not\prec A\Longrightarrow B\not\prec A';
\]

(iv) for each nonempty finite \(\{C_1,\ldots,C_m\}\) and every choice
of admitted profile continuations \(C_{-i}\) of
\(\{C_j:j\ne i\}\), at least one \(i\) satisfies
\(C_{-i}\not\prec C_i\).
If the required continuations do not exist, that choice makes no demand.
This is the finite-set reading: repeated names for the same argument
do not create additional members. An indexed-family reading instead
allows \(C_i=C_j\) for different indices, deleting only the selected
index when forming \(C_{-i}\). The countermodel below distinguishes
these readings; they must not be interchanged in a source theorem.
\end{hbdefinition}

\begin{hbdefinition}{def:p18-rationality-corrected-lifts}{The corrected strict preference convention}
Let \(<\) be the strict part of the appropriate premise or rule preorder.
For finite sets, define \(\triangleleft_{\rm E}\) and
\(\triangleleft_{\rm D}\) with these empty-set clauses: an empty left
operand is never below another operand; a nonempty operand is below
an empty right operand. For nonempty operands use
\[
\begin{aligned}
 \Gamma\triangleleft_{\rm E}\Delta
   &\iff(\exists x\in\Gamma)(\forall y\in\Delta)\ x<y,\\
 \Gamma\triangleleft_{\rm D}\Delta
   &\iff(\forall x\in\Gamma)(\exists y\in\Delta)\ x<y.
\end{aligned}
\]
The corrected last-link strict relation is
\[
 B\prec_{\rm L}^{\,c}A\iff
 \operatorname{Last}(B)\triangleleft\operatorname{Last}(A)
 \ \lor\
 \bigl(\operatorname{Last}(A)=\operatorname{Last}(B)=\varnothing
 \land\operatorname{Prem}_p(B)\triangleleft\operatorname{Prem}_p(A)\bigr).
\]
The corrected weakest-link strict relation uses, in order: premise
comparison if both arguments are strict; defeasible-rule comparison if
both are firm and not both strict; and the conjunction of the two strict
comparisons otherwise. These are direct strict definitions. They are
not the asymmetric parts of the tutorial's weak comparisons.
\end{hbdefinition}

\begin{hbtheorem}{thm:p18-rationality-strict-sets}{proved=annex:p18.ratsets}{Strict set comparisons and finite-union witnesses}
Both corrected set comparisons are strict partial orders.
If \(\bigcup_{i=1}^m\Gamma_i\triangleleft\Delta\), with \(m\geq1\),
then \(\Gamma_i\triangleleft\Delta\) for at least one \(i\).
For the democratic comparison every nonempty \(\Gamma_i\) works.
\end{hbtheorem}

\begin{hbtheorem}{thm:p18-rationality-transposition-boundary}{proved=annex:p18.ratcounter}{A boundary of the finite-set ordering criterion}
There is an unrestricted, finite structured framework with paired
contradictories, consistent axioms, strict rules closed under
transposition, and a strict partial preference satisfying all four
conditions of Definition~\ref{def:p18-rationality-reasonable}, whose
unique defeat-complete extension is inconsistent.
It is also defeat-grounded, defeat-preferred, and defeat-stable.
Its argument universe has 22 trees when all antecedent permutations
are represented explicitly. The proof supplies the full construction
and shows that its order fails the indexed-family condition allowing
repetitions.
\end{hbtheorem}

\section{A sufficient return condition and its general theorem}

\begin{hbdefinition}{def:p18-rationality-return}{Local return of an unsuccessful attack}
An admitted argument universe satisfies the \emph{local return condition}
if, whenever \(X T Y\) but not \(X D Y\), there exists an admitted
argument \(W\) such that
\[
 W D X,\qquad
 (\forall Z)\,[Z D W\Longrightarrow(Z D X\lor Z D Y)].
 \tag{R}
\]
All quantified arguments belong to the same admitted universe.
This condition concerns a failed attack and an explicit returning defeat;
it is not an assumption that accepted conclusions are consistent.
In a finite framework it is decidable from the complete attack and
defeat tables. If every attack succeeds, it holds vacuously.
\end{hbdefinition}

\begin{hblemma}{thm:p18-rationality-joint-defense}{proved=annex:p18.ratdefense}{Jointly defended arguments cannot attack}
Assume (R). If \(E\) is defeat-conflict-free and \(X,Y\) are both
defended by \(E\), then neither \(X T Y\) nor \(Y T X\).
This includes \(X=Y\). Defense is always calculated with \(D\).
\end{hblemma}

\begin{hblemma}{thm:p18-rationality-support-union}{proved=annex:p18.ratsupport}{Consistent unions in the restricted universe}
Assume (R), c-classicality, and axiom consistency in a c-restricted
framework. If finitely many \(A_i\) are defended by a defeat-conflict-free
set \(E\), then \(\bigcup_i\operatorname{Prem}(A_i)\) is c-consistent.
\end{hblemma}

\begin{hbtheorem}{thm:p18-rationality-general}{proved=annex:p18.ratgeneral}{General closure, consistency, and semantic equivalence}
Use either all finite arguments or all c-consistent finite arguments.
Assume axiom consistency, contrary well-formedness, and (R).
Assume additionally either closure under contraposition or, for a
paired-literal language, closure under transposition. In the c-restricted
case also assume c-classicality.

Then attack-admissible and defeat-admissible extensions coincide, as do
complete, grounded, preferred, and stable extensions under the two
conflict-freeness conventions. Admissibility here means the corresponding
conflict-freeness together with defense of every member using \(D\).
Every admissible extension has directly consistent conclusions.
Every complete extension is closed under subarguments and strict rules,
and its conclusions are indirectly consistent.
The theorem applies to infinite argument sets; only individual trees and
strict-rule antecedents must be finite.
\end{hbtheorem}

The local return condition isolates the precise defeat-producing
requirement needed by the source's failed-attack argument.
The preceding theorem supplies a full replacement at that explicit
interface. The reasonable-order conditions are not silently promoted
to a proof of (R).
The axiom-consistency and strict-rule requirements also have substantive
content: the first two systems in
Theorem~\ref{thm:p18-aspic-rationality-boundaries} have every attack
succeed, so satisfy (R), but exhibit the corresponding consistency
failures when those requirements are absent. The third system there
explains why the present theorem keeps rebuttal confined to fallible
roots.

\section{The corrigendum's four-premise system}

Let \(L\) consist of paired literals over \(p,q,r,s\), let
\(K_p=\{p,q,r,s\}\), and let \(K_n=R_d=\varnothing\).
For each \(x\in\{p,q,r,s\}\), include the strict rule whose antecedents
are the other three positive literals and whose consequent is \(\neg x\).
Include all six permutations of each antecedent list.
The ordinary-premise preorder has two incomparable equivalence classes,
\(\{p,q\}\) and \(\{r,s\}\). Thus its strict part is empty.

\begin{hbtheorem}{thm:p18-rationality-corrigendum-example}{proved=annex:p18.ratexample}{All arguments and all extensions of the corrected example}
There are four premise arguments \(P_x=[x]\) and 24 rule arguments
\(N_{x,\sigma}\), one for each missing literal \(x\) and each permutation
\(\sigma\) of the other three literals.
Under either corrected link principle and either corrected set comparison,
the strict argument preference is empty and \(D=T\).
For each \(x\), put
\[
 E_x=\{P_y:y\ne x\}\cup\{N_{x,\sigma}:\sigma\text{ a permutation}\}.
\]
The four \(E_x\) are exactly the preferred and stable extensions.
The complete extensions are these four sets together with the empty set;
the grounded extension is empty.
The unrestricted and c-restricted argument universes are identical here,
and all hypotheses of Theorem~\ref{thm:p18-rationality-general} hold.

With the tutorial's elitist weak comparison followed by its asymmetric
part, every attack fails instead. The entire 28-argument universe is then
the unique defeat-complete and defeat-stable extension, and its conclusions
are inconsistent. Thus the two source preference conventions cannot be
interchanged.
\end{hbtheorem}

\begin{proofboundary}{def:p18-rationality-next}{Deriving rationality from specified priorities}
The finite-set reading of the corrected source's ordering assumptions
does not support the transposition branch, by
Theorem~\ref{thm:p18-rationality-transposition-boundary}.
The theorem established here uses the explicit additional condition (R).
For all four corrected priority constructions, the required implication
is supplied in Chapter~\ref{sec:p18-structured-priority-rationality}.
Theorems~\ref{thm:p18-priority-return} and
\ref{thm:p18-priority-rationality} establish global return for
unrestricted trees and the support-qualified route for c-restricted
trees, under each stated strict-rule regime.

The occurrence-preserving reversal is proved in
Theorem~\ref{thm:p18-priority-retained-reversal}; it retains the
assumptions needed for actual grafts. The literal profile-continuation
condition remains distinct from actual strict tree continuation.
Classical/Tarskian representation theorems and preferred-subtheory
equivalence are separate source presentations. The direct set-comparison
theorem and the complete corrigendum example above retain their own
stated domains.
\end{proofboundary}
